\part{Yarncrawler: Self-Repair as Universal Principle}
\label{part:yarncrawler}

\epigraph{To act intelligently is to Yarncrawl: to repair oneself by weaving meaning
into structure, to sustain homeorhetic flows rather than static equilibria, and to
treat skepticism not as a threat but as a guide to resilience.}{}

\noindent Part~\ref{part:yarncrawler} develops the Yarncrawler framework as the
self-repair layer of the entire theoretical architecture. The Yarncrawler is not an
application of the foregoing concepts but the mechanism by which they maintain
coherence with each other. The sheaf-theoretic foundation of Chapter~\ref{ch:sheaf}
provides the formal statement of what the monograph is doing to itself: stitching
overlapping coordinate charts into a global section without destroying the generative
richness of the local descriptions.
\cleardoublepage
